% ---- V. Experimental Setup. Budget 1.2 pages. Source: JAMA "Study Design", "Simulation Model
% Overview", "Action Space" (hospital verification sentence). Mostly NEW detail the reviewers
% asked for: cohort characteristics, replicates, hospital verification, sensitivity design.
\section{Experimental Setup}
\label{sec:setup}

\subsection{Regions and hospital networks}
\jama{JAMA "Study Design and Setting" + hospital-verification sentence, updated to the 2 Oct 2026 verification.}
We evaluate two regions with different network structures and stroke-system policies: the nine-county San Francisco Bay Area, California, a decentralised system with \pending{11} EVT-capable centres, \pending{44} thrombolysis-capable centres and \pending{6} non-stroke-centre hospitals (\pending{61} in total), and the state of Rhode Island, a centralised system with one EVT-capable centre, eight thrombolysis-capable centres and no uncertified acute-care hospital. \new{Hospital designations were verified on 2 October 2026 against county EMS agency receiving-facility lists, the Rhode Island Department of Health stroke designation list, certifying-body records and hospital system pages; coordinates were checked against street addresses. Four Bay Area hospitals perform thrombectomy without Comprehensive or Thrombectomy-Capable certification; they are thrombolysis-capable in the primary analysis and EVT-capable in a sensitivity analysis. Designations change (two Bay Area hospitals changed in 2026), so the hospital table is an input to the method rather than part of it.} \new{[Add one sentence on EMS context: SF routes to the closest stroke centre without LVO tiering; Santa Clara routes LAMS-positive patients to the closest EVT-capable centre; Marin transfers LVO patients out of county, often by air. Cite policies.]}

\subsection{Simulated cohorts}
\jama{JAMA "Simulation Model Overview", with replicate detail added (Reviewer 2: how many simulations, what characteristics).}
Each replicate samples 1,000 pickup locations without replacement from a population-weighted pool of 10,000 points~\cite{meta2019}, draws onset-to-pickup time uniformly on 30--270 minutes and the subtype from the LAMS-positive prior, and routes every patient under each policy. \new{We run 10 replicates per region with distinct seeds, giving 10,000 patient-policy pairs per policy and between-replicate confidence intervals. The same patients are routed under every policy, so comparisons are paired. A patient is excluded only if no policy can reach a hospital within the catchment radius (\pending{k} of 10,000).} \new{[Reviewer 1 asked about the grid: state here that the 200$\times$200 grid with 5 patients per cell is the visualisation layer for the geographic analysis, not the cohort used for the main results.]}

\subsection{Travel times and catchment}
\new{Travel times are computed on the OpenStreetMap road network with OpenRouteService~8.0 (driving-car profile) hosted locally. Candidate destinations for the initial transport are restricted to facilities within 80~km of the pickup location, an EMS catchment assumption; inter-facility transfers are computed by road for all pairs without a distance restriction. Travel times are deterministic in the planner; sensitivity to travel-time error is assessed post hoc (below).}

\subsection{Outcome measures and statistics}
\new{The primary measure is the probability of an excellent 90-day outcome (mRS 0--1) per patient. For each pair of policies we report the mean paired difference with a 95\% confidence interval across the 10 replicates ($t_9$), the paired $t$-test with Bonferroni correction for the five comparisons, and Cohen's $d_z$. We distinguish the \emph{planning} reward (the prior-weighted mixture the planner optimises) from the \emph{realised} reward (the subtype-specific outcome given the drawn subtype); tables report realised rewards unless stated. Equity is assessed by stratifying paired differences by census-tract median household income quartile and RUCA urbanicity, for the MDP against nearest-hospital and against both heuristics.}

\subsection{Sensitivity analyses}
\new{(i) Travel-time error: multiplicative log-normal noise ($\sigma = 0.10, 0.20$) and systematic bias ($\times 0.70$ to $\times 1.30$) applied to the recorded first-leg travel time, with one draw per patient-destination so identical routes receive identical perturbations; rewards are recomputed and the paired comparisons repeated. (ii) In-hospital intervals: $T_{\text{DIDO}} \in \{60, 90, 121, 175\}$ minutes, rewards recomputed post hoc for the policies as planned under the base case. (iii) EVT definition: the four uncertified thrombectomy sites treated as EVT-capable, with the policy re-planned. (iv) Decision-tree depth: 1 to 10 and unconstrained.}

\subsection{Latency}
\new{Per-decision wall-clock time of the planner is measured on \pending{machine spec} for 50 patients at planning depths 1 to 4, with the routing engine on the same machine, after warm-up. We report mean, median, 95th and 99th percentiles.}

\new{[Reproducibility sentence: all randomness seeded; pinned Julia environment; code and hospital tables released on acceptance.]}
